% GENERATED FILE. Edit canonical lessons, then run npm run generate:books.
% canonical-source-hash: fnv1a64:d7176d874abc7693
% canonical-lessons: MR-C212-arj, MR-C212-pakit, MR-C212-patta, MR-C212-pagar, MR-C212-margadarsak

\chapter{Application, Envelope, Address, Salary, Guide}
\label{ch:mr-application-envelope-address-salary-guide}
\hlchaptermodality{212}

\begin{chapteropening}
\textbf{By the end of this chapter:} \emph{I can say an application, an envelope, an address, a salary and a guide.}

Complete the last lesson of chapter 212: I can say an application, an envelope, an address, a salary and a guide.
\end{chapteropening}

\section[\texorpdfstring{\mr{अर्ज} (arj)}{अर्ज (arj)}]{\mr{अर्ज} (arj) — an application, a form}

\label{lesson:MR-C212-arj}



\begin{quote}
Before the new one: say the Marathi for a decision, then the Marathi for a possibility.
\end{quote}

\subsection*{You'll want to know: \mr{अर्ज}}
\textbf{\mr{अर्ज}} — \emph{arj} — \textquotedblleft{}an application, a form\textquotedblright{}.

\begin{cousinweb}
\textbf{\mr{अर्ज}} (\emph{arj}) is a loan from Arabic \emph{ʿarḍ}, \textquotedblleft{}a presenting, a petition\textquotedblright{}, which reached Marathi through Persian. Its Arabic root is \emph{ʿ-r-ḍ} (Lane). Persian says \emph{ʿarz}.
\end{cousinweb}

\begin{grammarlens}[title={What you've built}]
One more word for this chapter.
\end{grammarlens}

\subsection*{Your turn}
\begin{itemize}
\raggedright
  \item \emph{Say it:} \emph{arj}
  \item \emph{Say it:} \emph{arj}, once more
  \item \emph{Say it:} say \emph{arj}
  \item \emph{Recall:} say the Marathi for a decision, then the Marathi for a possibility, then say \emph{arj} again
\end{itemize}

\begin{tcolorbox}[breakable,colback=teal!4,colframe=teal!35!black,title={Before you move on}]
What does \mr{अर्ज} mean? (\textquotedblleft{}An application, a form\textquotedblright{}.) Say it once more.
\end{tcolorbox}
\section[\texorpdfstring{\mr{पाकीट} (pākīṭ)}{पाकीट (pākīṭ)}]{\mr{पाकीट} (pākīṭ) — an envelope; a wallet}

\label{lesson:MR-C212-pakit}



\begin{quote}
Before the new one: say the Marathi for a possibility, then the Marathi for an application.
\end{quote}

\subsection*{You'll want to know: \mr{पाकीट}}
\textbf{\mr{पाकीट}} — \emph{pākīṭ} — \textquotedblleft{}an envelope; a wallet\textquotedblright{}.

\begin{cousinweb}
\textbf{\mr{पाकीट}} (\emph{pākīṭ}) is a loan from English \emph{packet}.
\end{cousinweb}

\begin{grammarlens}[title={What you've built}]
One more word for this chapter.
\end{grammarlens}

\subsection*{Your turn}
\begin{itemize}
\raggedright
  \item \emph{Say it:} \emph{pākīṭ}
  \item \emph{Say it:} \emph{pākīṭ}, once more
  \item \emph{Say it:} say \emph{pākīṭ}
  \item \emph{Recall:} say the Marathi for a possibility, then the Marathi for an application, then say \emph{pākīṭ} again
\end{itemize}

\begin{tcolorbox}[breakable,colback=teal!4,colframe=teal!35!black,title={Before you move on}]
What does \mr{पाकीट} mean? (\textquotedblleft{}An envelope; a wallet\textquotedblright{}.) Say it once more.
\end{tcolorbox}
\section[\texorpdfstring{\mr{पत्ता} (pattā)}{पत्ता (pattā)}]{\mr{पत्ता} (pattā) — an address}

\label{lesson:MR-C212-patta}



\begin{quote}
Before the new one: say the Marathi for an application, then the Marathi for an envelope; a wallet.
\end{quote}

\subsection*{You'll want to know: \mr{पत्ता}}
\textbf{\mr{पत्ता}} — \emph{pattā} — \textquotedblleft{}an address\textquotedblright{}.

\begin{cousinweb}
The origin of \textbf{\mr{पत्ता}} (\emph{pattā}), \textquotedblleft{}address\textquotedblright{}, is not settled. CDIAL has no entry for it in this sense.
\end{cousinweb}

\begin{grammarlens}[title={What you've built}]
One more word for this chapter.
\end{grammarlens}

\subsection*{Your turn}
\begin{itemize}
\raggedright
  \item \emph{Say it:} \emph{pattā}
  \item \emph{Say it:} \emph{pattā}, once more
  \item \emph{Say it:} say \emph{pattā}
  \item \emph{Recall:} say the Marathi for an application, then the Marathi for an envelope; a wallet, then say \emph{pattā} again
\end{itemize}

\begin{tcolorbox}[breakable,colback=teal!4,colframe=teal!35!black,title={Before you move on}]
What does \mr{पत्ता} mean? (\textquotedblleft{}An address\textquotedblright{}.) Say it once more.
\end{tcolorbox}
\section[\texorpdfstring{\mr{पगार} (pagār)}{पगार (pagār)}]{\mr{पगार} (pagār) — a salary}

\label{lesson:MR-C212-pagar}



\begin{quote}
Before the new one: say the Marathi for an envelope; a wallet, then the Marathi for an address.
\end{quote}

\subsection*{You'll want to know: \mr{पगार}}
\textbf{\mr{पगार}} — \emph{pagār} — \textquotedblleft{}a salary\textquotedblright{}.

\begin{cousinweb}
\textbf{\mr{पगार}} (\emph{pagār}) is a loan from Portuguese \emph{pagar}, \textquotedblleft{}to pay\textquotedblright{}.
\end{cousinweb}

\begin{grammarlens}[title={What you've built}]
One more word for this chapter.
\end{grammarlens}

\subsection*{Your turn}
\begin{itemize}
\raggedright
  \item \emph{Say it:} \emph{pagār}
  \item \emph{Say it:} \emph{pagār}, once more
  \item \emph{Say it:} say \emph{pagār}
  \item \emph{Recall:} say the Marathi for an envelope; a wallet, then the Marathi for an address, then say \emph{pagār} again
\end{itemize}

\begin{tcolorbox}[breakable,colback=teal!4,colframe=teal!35!black,title={Before you move on}]
What does \mr{पगार} mean? (\textquotedblleft{}A salary\textquotedblright{}.) Say it once more.
\end{tcolorbox}
\section[\texorpdfstring{\mr{मार्गदर्शक} (mārgadarśak)}{मार्गदर्शक (mārgadarśak)}]{\mr{मार्गदर्शक} (mārgadarśak) — a guide}

\label{lesson:MR-C212-margadarsak}



\begin{quote}
Before the new one: say the Marathi for an address, then the Marathi for a salary.
\end{quote}

\subsection*{You'll want to know: \mr{मार्गदर्शक}}
\textbf{\mr{मार्गदर्शक}} — \emph{mārgadarśak} — \textquotedblleft{}a guide\textquotedblright{}.

\begin{cousinweb}
\textbf{\mr{मार्गदर्शक}} (\emph{mārgadarśak}) is a compound of Sanskrit parts: \emph{mārga}, \textquotedblleft{}track, road\textquotedblright{} (CDIAL 10071), and \emph{darśaka}, \textquotedblleft{}one who shows\textquotedblright{}, from \emph{dṛś}, \textquotedblleft{}to see\textquotedblright{}.
\end{cousinweb}

\begin{grammarlens}[title={What you've built}]
One more word for this chapter.
\end{grammarlens}

\subsection*{Your turn}
\begin{itemize}
\raggedright
  \item \emph{Say it:} \emph{mārgadarśak}
  \item \emph{Say it:} \emph{mārgadarśak}, once more
  \item \emph{Say it:} say \emph{mārgadarśak}
  \item \emph{Recall:} say the Marathi for an address, then the Marathi for a salary, then say \emph{mārgadarśak} again
\end{itemize}

\begin{tcolorbox}[breakable,colback=teal!4,colframe=teal!35!black,title={Before you move on}]
What does \mr{मार्गदर्शक} mean? (\textquotedblleft{}A guide\textquotedblright{}.) Say it once more.
\end{tcolorbox}
